% Emoji-settings.tex - Moderne Emoji-Unterstützung
% Ergänzung zum modernisierten LaTeX-Template

% === Emoji-Font-System ===
% Automatische Erkennung verfügbarer Emoji-Fonts
\newif\ifEmojiFontFound
\EmojiFontFoundfalse

\IfFontExistsTF{Noto Color Emoji}{
  \newfontfamily{\emojifont}{Noto Color Emoji}[
    Renderer=HarfBuzz,
    Scale=MatchUppercase
  ]
  \typeout{Emoji-Font: Noto Color Emoji}
  \EmojiFontFoundtrue
}{
  \IfFontExistsTF{Apple Color Emoji}{
    \newfontfamily{\emojifont}{Apple Color Emoji}[
      Renderer=HarfBuzz,
      Scale=MatchUppercase
    ]
    \typeout{Emoji-Font: Apple Color Emoji}
    \EmojiFontFoundtrue
  }{
    \IfFontExistsTF{Segoe UI Emoji}{
      \newfontfamily{\emojifont}{Segoe UI Emoji}[
        Renderer=HarfBuzz,
        Scale=MatchUppercase
      ]
      \typeout{Emoji-Font: Segoe UI Emoji}
      \EmojiFontFoundtrue
    }{
      \IfFontExistsTF{Twemoji}{
        \newfontfamily{\emojifont}{Twemoji}[
          Renderer=HarfBuzz,
          Scale=MatchUppercase
        ]
        \typeout{Emoji-Font: Twemoji}
        \EmojiFontFoundtrue
      }{
        % Fallback: Verwende Hauptfont
        \let\emojifont\rmfamily
        \PackageWarning{emoji-settings}{%
          Kein Emoji-Font gefunden. Emojis werden möglicherweise nicht korrekt dargestellt.
        }
        \typeout{Emoji-Font: Fallback (Hauptfont, da keine Emoji-Font gefunden)}
      }
    }
  }
}

% === Robustes Success-Emoji mit TikZ-Fallback ===
\makeatletter
\@ifpackageloaded{tikz}{
  \newcommand{\drawSuccess}{%
    \begin{tikzpicture}[baseline=-0.5ex, scale=0.25]
      \fill[success] (0,0) circle (1.2);
      \draw[white, line width=0.5mm] (-0.5,0) -- (-0.2,-0.4) -- (0.6,0.5);
    \end{tikzpicture}%
  }
  \newcommand{\drawFailure}{%
     \begin{tikzpicture}[baseline=-0.5ex, scale=0.25]
      \fill[danger] (0,0) circle (1.2);
      \draw[white, line width=0.5mm] (-0.4,-0.4) -- (0.4,0.4);
      \draw[white, line width=0.5mm] (-0.4,0.4) -- (0.4,-0.4);
    \end{tikzpicture}%
  }
  \newcommand{\drawWarning}{%
    \begin{tikzpicture}[baseline=-0.3ex, scale=0.25]
      \node[regular polygon, regular polygon sides=3, draw=warning, fill=warning, inner sep=0.5pt, rounded corners=0.5pt] at (0,0) {};
      \node[black, font=\bfseries\tiny] at (0,-0.1) {!};
    \end{tikzpicture}%
  }
  \newcommand{\drawInfo}{%
    \begin{tikzpicture}[baseline=-0.5ex, scale=0.25]
      \fill[info] (0,0) circle (1.2);
      \node[white, font=\bfseries\large] at (0,0) {i};
    \end{tikzpicture}%
  }
}{
  \newcommand{\drawSuccess}{\textcolor{success}{Ok}}
  \newcommand{\drawFailure}{\textcolor{danger}{X}}
  \newcommand{\drawWarning}{\textcolor{warning}{!}}
  \newcommand{\drawInfo}{\textcolor{info}{i}}
}
\makeatother

% === Basis-Emoji-Befehl ===
\newcommand{\emoji}[1]{{\emojifont #1}}

% === Häufige Emojis ===

% Definiere \success und Co. robust
\ifEmojiFontFound
  \providecommand{\success}{\textcolor{success}{\emoji{✅}}}
  \providecommand{\failure}{\textcolor{danger}{\emoji{❌}}}
  \providecommand{\warning}{\textcolor{warning}{\emoji{⚠️}}}
  \providecommand{\info}{\textcolor{info}{\emoji{ℹ️}}}
\else
  % Wenn kein Emoji-Font gefunden, nutze TikZ
  \providecommand{\success}{\drawSuccess}
  \providecommand{\failure}{\drawFailure}
  \providecommand{\warning}{\drawWarning}
  \providecommand{\info}{\drawInfo}
\fi


% Daumen & Gesten (nur wenn Font da)
\providecommand{\thumbsup}{\ifEmojiFontFound\emoji{}\else(ok)\fi}
\providecommand{\thumbsdown}{\ifEmojiFontFound\emoji{}\else(no)\fi}
\providecommand{\clap}{\ifEmojiFontFound\emoji{}\else(clap)\fi}
\providecommand{\wave}{\ifEmojiFontFound\emoji{}\else(wave)\fi}
\providecommand{\point}{\ifEmojiFontFound\emoji{}\else(point)\fi}

% Herzen & Sterne (Prefix \e...)
\providecommand{\heart}{\ifEmojiFontFound\emoji{❤️}\else<3\fi}
\newcommand{\eHeartRed}{\textcolor{danger}{\ifEmojiFontFound\emoji{❤️}\else<3\fi}}
\newcommand{\eStar}{\ifEmojiFontFound\emoji{⭐}\else*\fi} 
\newcommand{\eStarFull}{}

% Tech & Business
\providecommand{\rocket}{}
\providecommand{\fire}{}
\providecommand{\lightning}{\ifEmojiFontFound\emoji{⚡}\else flash \fi}
\providecommand{\gear}{\ifEmojiFontFound\emoji{⚙️}\else gear \fi}
\providecommand{\computer}{\ifEmojiFontFound\emoji{}\else PC \fi}
\providecommand{\mobile}{\ifEmojiFontFound\emoji{}\else Handy \fi}

% Semantische Status Alternative (Texte)
\providecommand{\error}{\failure}
\providecommand{\question}{\textcolor{secondary}{?}}
\providecommand{\attention}{\warning}

% Alternative Darstellungen
\providecommand{\tick}{\success}
\providecommand{\checkmark}{\success}
\providecommand{\crossmark}{\failure}

% === Fachspezifische Symbole ===
% Cloud & IT
\providecommand{\cloud}{\ifEmojiFontFound\emoji{☁️}\else Cloud \fi}
\providecommand{\database}{\ifEmojiFontFound\emoji{️}\else DB \fi}
\providecommand{\server}{\ifEmojiFontFound\emoji{️}\else Server \fi}
\providecommand{\network}{\ifEmojiFontFound\emoji{}\else Network \fi}
\providecommand{\security}{\ifEmojiFontFound\emoji{}\else Security \fi}
\providecommand{\key}{\ifEmojiFontFound\emoji{}\else Key \fi}

% Business & Kommunikation
\providecommand{\email}{\ifEmojiFontFound\emoji{}\else Email \fi}
\providecommand{\phone}{\ifEmojiFontFound\emoji{}\else Phone \fi}
\providecommand{\calendar}{\ifEmojiFontFound\emoji{}\else Cal \fi}
\providecommand{\clock}{\ifEmojiFontFound\emoji{}\else Time \fi}
\providecommand{\location}{\ifEmojiFontFound\emoji{}\else Loc \fi}
\providecommand{\building}{\ifEmojiFontFound\emoji{}\else Bldg \fi}

% Dokumente & Dateien
\newcommand{\eDocument}{}
\newcommand{\eFolder}{}
\newcommand{\eChart}{}
\newcommand{\eGraph}{}
\newcommand{\eTrending}{}


% === Pfeile & Navigation ===
\newcommand{\eArrowRight}{\ifEmojiFontFound\emoji{→}\else $\rightarrow$ \fi}
\newcommand{\eArrowLeft}{\ifEmojiFontFound\emoji{←}\else $\leftarrow$ \fi}
\newcommand{\eArrowUp}{\ifEmojiFontFound\emoji{↑}\else $\uparrow$ \fi}
\newcommand{\eArrowDown}{\ifEmojiFontFound\emoji{↓}\else $\downarrow$ \fi}
\newcommand{\eArrowLeftRight}{\ifEmojiFontFound\emoji{↔}\else $\leftrightarrow$ \fi}
\newcommand{\eImplies}{\ifEmojiFontFound\emoji{⇒}\else $\Rightarrow$ \fi}

% Farbige Pfeile
\newcommand{\eGreenArrow}{\textcolor{success}{\eArrowRight}}
\newcommand{\eRedArrow}{\textcolor{danger}{\eArrowRight}}
\newcommand{\eBlueArrow}{\textcolor{primary}{\eArrowRight}}

% === Listen mit Emoji-Bullets ===
\newlist{emojilist}{itemize}{3}
\setlist[emojilist,1]{label=, leftmargin=2em}
\setlist[emojilist,2]{label=▫️, leftmargin=1.5em}
\setlist[emojilist,3]{label=•, leftmargin=1em}

\newlist{checklist}{itemize}{2}
\setlist[checklist,1]{label=\success, leftmargin=2em}
\setlist[checklist,2]{label=\checkmark, leftmargin=1.5em}

\newlist{todolist}{itemize}{2}
\setlist[todolist,1]{label=☐, leftmargin=2em}
\setlist[todolist,2]{label=▫️, leftmargin=1.5em}

\newlist{starlist}{itemize}{2}
\setlist[starlist,1]{label=\eStar, leftmargin=2em}
\setlist[starlist,2]{label=, leftmargin=1.5em}

% === Spezielle Formate ===
\newcommand{\important}[1]{%
  \warning \textbf{\textcolor{warning}{Wichtig:}} #1%
}

\newcommand{\note}[1]{%
  \info \textbf{\textcolor{info}{Hinweis:}} #1%
}

\newcommand{\tip}[1]{%
   \textbf{\textcolor{primary}{Tipp:}} #1%
}

\newcommand{\achievement}[1]{%
   \textbf{\textcolor{success}{Erfolg:}} #1%
}

% === Bewertungs-Systeme ===
\newcommand{\onestar}{\eStar}
\newcommand{\twostars}{\eStar\eStar}
\newcommand{\threestars}{\eStar\eStar\eStar}
\newcommand{\fourstars}{\eStar\eStar\eStar\eStar}
\newcommand{\fivestars}{\eStar\eStar\eStar\eStar\eStar}

% Ampel-System
\newcommand{\redlight}{\textcolor{danger}{\ifEmojiFontFound\emoji{}\else(ROT)\fi}}
\newcommand{\yellowlight}{\textcolor{warning}{\ifEmojiFontFound\emoji{}\else(GELB)\fi}}
\newcommand{\greenlight}{\textcolor{success}{\ifEmojiFontFound\emoji{}\else(GRUEN)\fi}}

% === Test-Funktion ===
\newcommand{\emojitest}{%
  \begin{center}
  \textbf{Emoji-System Test (TikZ Fallback: \ifEmojiFontFound nein \else ja \fi)}\\[0.5em]
  \begin{tabular}{ll}
  \textbf{Status:} & \success{} \failure{} \warning{} \info{} \\
  \end{tabular}
  \end{center}
}

% === Beispiel-Umgebungen ===
\newenvironment{successbox}{%
  \begin{tcolorbox}[
    enhanced,
    colback=success!10,
    colframe=success,
    boxrule=1pt,
    sharp corners,
    before upper={\success{} \textbf{Erfolg:} }
  ]
}{%
  \end{tcolorbox}
}

\newenvironment{warningbox}{%
  \begin{tcolorbox}[
    enhanced,
    colback=warning!10,
    colframe=warning,
    boxrule=1pt,
    sharp corners,
    before upper={\warning{} \textbf{Warnung:} }
  ]
}{%
  \end{tcolorbox}
}

\newenvironment{infobox}{%
  \begin{tcolorbox}[
    enhanced,
    colback=info!10,
    colframe=info,
    boxrule=1pt,
    sharp corners,
    before upper={\info{} \textbf{Information:} }
  ]
}{%
  \end{tcolorbox}
}

% === Integration mit bestehendem System ===
\makeatletter
\@ifundefined{success}{
  \definecolor{success}{RGB}{40, 167, 69}
  \definecolor{danger}{RGB}{220, 53, 69}
  \definecolor{info}{RGB}{23, 162, 184}
  \definecolor{warning}{RGB}{255, 193, 7}
  \definecolor{primary}{RGB}{0, 102, 153}
  \definecolor{secondary}{RGB}{102, 102, 102}
}{}
\makeatother

% === Debug-Information ===
\AtBeginDocument{%
  \typeout{=== Emoji-System aktiviert ===}
  \ifEmojiFontFound
     \typeout{Emoji-Font: Gefunden (Unicode)}
  \else
     \typeout{Emoji-Font: NICHT GEFUNDEN (TikZ Fallback aktiv)}
  \fi
  \typeout{===============================}
}